\section{Introducción}
El análisis de ratios financieros permite convertir cifras contables en indicadores útiles para evaluar una empresa. En el contexto de Contabilidad Empresarial, los ratios ayudan a responder preguntas concretas: si la entidad puede pagar sus obligaciones corrientes, si cobra oportunamente, si depende demasiado de deuda y si sus recursos generan utilidad.

Este documento desarrolla los principales ratios de liquidez, gestión, solvencia y rentabilidad. La teoría se presenta de manera breve y se aplica inmediatamente en casos académicos contextualizados en empresas peruanas. Todos los casos son ficticios y han sido elaborados con fines educativos; por tanto, no representan estados financieros reales.

El enfoque del trabajo no consiste solo en calcular. Cada resultado se acompaña de interpretación y decisión empresarial, porque un ratio aislado puede llevar a conclusiones incompletas. Un indicador debe contrastarse con políticas internas, periodos anteriores, metas, sector económico, calidad de la información contable y capacidad real de generar efectivo.
